\chapter{Soal Cerita dan Diagram}\label{ch:g5:data}

Bab terakhir tahun ini memakai semuanya: soal cerita bertingkat
tentang uang, banyaknya barang dan lama waktu, serta diagram yang
menyajikan keterangan dalam sekali pandang. Gagasan utamanya adalah
``per satu'' --- mencari harga satu barang --- yang kelak tumbuh
menjadi perbandingan senilai pada \cref{ch:g6:propdata}.

\section{Soal cerita bertingkat}

\begin{method}[Menjinakkan soal yang panjang]\label{met:g5:data:steps}
\begin{enumerate}
  \item Bacalah sampai habis; katakan apa pertanyaan akhirnya;
  \item rencanakan mundur: untuk menjawabnya, apa yang kuperlukan
        lebih dulu?
  \item selesaikan langkahnya berurutan, dengan satu operasi dan
        satu kalimat pendek untuk tiap langkah;
  \item periksalah jawaban akhirnya dengan akal sehat dan dengan
        taksiran.
\end{enumerate}
\end{method}

\begin{example}\label{ex:g5:data:multistep}
``Sebuah sekolah memesan $4$ kotak berisi $25$ buku seharga $6$
masing-masing, dan membayar $180$ lagi untuk pengiriman. Berapa
seluruh biayanya?''
\begin{enumerate}
  \item Buku: $4 \times 25 = 100$ buku.
  \item Biaya buku: $100 \times 6 = 600$.
  \item Seluruhnya: $600 + 180 = 780$.
\end{enumerate}
Periksa dengan taksiran: kira-kira seratus buku seharga $6$ berarti
kira-kira $600$, ditambah kira-kira $200$: sekitar $800$ --- jawaban $780$ masuk akal.
\end{example}

\begin{example}[Per satu]\label{ex:g5:data:perone}
``$5$ cangkir yang sama berharga $12.50$; berapa harga $8$ cangkir?''
Carilah dulu harga \emph{satu} cangkir:
\[
12.50 \div 5 = 2.50,
\qquad\text{lalu}\qquad
8 \times 2.50 = 20 .
\]
Lewat satuan, semuanya menjadi terjangkau --- pola dua langkah ini
menyelesaikan satu keluarga soal (resep, kelajuan, harga), dan
menjadi benih \cref{ch:g6:propdata}.
\end{example}

\begin{example}[Kembalian dan anggaran]\label{ex:g5:data:budget}
Zahra punya $30$. Ia membeli buku seharga $12.40$ dan dua pena
seharga $2.30$ masing-masing. Sanggupkah ia membeli permainan papan seharga $13$ juga?
\begin{enumerate}
  \item Pena: $2 \times 2.30 = 4.60$.
  \item Sudah dibelanjakan: $12.40 + 4.60 = 17$.
  \item Sisanya: $30 - 17 = 13$ --- pas cukup untuk permainan itu,
        tanpa \omterm{def:g3:division:remainder}{sisa} sama sekali.
\end{enumerate}
\end{example}

\section{Membaca diagram}

\begin{method}[Membaca diagram apa pun]\label{met:g5:data:read}
\begin{enumerate}
  \item Bacalah judulnya: apa yang sedang dibilang?
  \item Bacalah sumbunya atau kuncinya: satu garis skala, satu
        batang, satu lambang mewakili apa?
  \item Barulah menjawab pertanyaannya --- menunjuk diagramnya
        dengan penggaris membantu membaca tingginya dengan tepat.
\end{enumerate}
\end{method}

\begin{example}\label{ex:g5:data:chart}
Curah hujan di sebuah kota, bulan demi bulan:

\begin{omfigure}
\begin{tikzpicture}
\begin{axis}[
  width=8.8cm, height=5.0cm,
  ybar, bar width=13pt,
  symbolic x coords={Jan,Feb,Mar,Apr,Mei,Jun},
  xtick=data,
  ymin=0, ymax=90,
  ytick={20,40,60,80},
  ylabel={hujan (mm)},
  tick label style={font=\small},
  label style={font=\small}]
  \addplot[fill=omDef!30, draw=omDef] coordinates
    {(Jan,60) (Feb,45) (Mar,55) (Apr,70) (Mei,80) (Jun,30)};
\end{axis}
\end{tikzpicture}

{\small Diagram batang curah hujan. Satu garis skala $= 20$ mm.}
\end{omfigure}

Bacaan: bulan paling basah adalah Mei ($80$ mm); Juni mendapat
setengah hujan April ($30$ melawan $70$? bukan --- setengah dari
$70$ adalah $35$, jadi \emph{tidak persis} setengah: membaca dengan
cermat itu penting!); seluruhnya untuk enam bulan:
$60 + 45 + 55 + 70 + 80 + 30 = 340$ mm.
\end{example}

\begin{example}[Besaran yang berubah]\label{ex:g5:data:line}
Tinggi sebuah tanaman, diukur setiap minggu:

\begin{omfigure}
\begin{tikzpicture}
\begin{axis}[omaxis,
  width=8.4cm, height=5.0cm,
  xmin=0, xmax=6.5, ymin=0, ymax=24,
  xlabel={minggu}, ylabel={tinggi (cm)},
  xtick={1,...,6}, ytick={5,10,15,20}]
  \addplot[omDef, thick, mark=*, mark size=1.6pt] coordinates
    {(1,3) (2,6) (3,10) (4,15) (5,19) (6,21)};
\end{axis}
\end{tikzpicture}

{\small Menghubungkan titik pengukurannya memperlihatkan
\emph{pertumbuhannya}: makin curam ruas garisnya, makin cepat tanaman itu tumbuh pada minggu tersebut.}
\end{omfigure}

Tanaman itu tumbuh paling cepat antara minggu $3$ dan $4$ ($+5$ cm)
dan melambat di akhir ($+2$ cm pada minggu $6$).
\end{example}

\section{Latihan}

\begin{exercise}[$\star$]\label{exo:g5:data:1}
Selesaikan dengan langkah: ``Sebuah klub menyewa bus seharga $240$
dan membeli $32$ karcis museum seharga $7$ masing-masing. Berapa seluruh biaya perjalanannya?''
\end{exercise}

\begin{exercise}[$\star$]\label{exo:g5:data:2}
Selesaikan dengan cara ``per satu'': ``$3$ kg apel berharga $7.50$.
Berapa harga $5$ kg?''
\end{exercise}

\begin{exercise}[$\star$]\label{exo:g5:data:3}
Selesaikan: ``Biaya piknik sebesar $87$ dibagi rata kepada enam
orang teman. Berapa yang dibayar masing-masing?''
\end{exercise}

\begin{exercise}[$\star$]\label{exo:g5:data:4}
Sandi punya $25$. Ia membeli kaos seharga $13.90$ dan topi seharga
$8.50$. Berapa sisa uangnya?
\end{exercise}

\begin{exercise}[$\star$]\label{exo:g5:data:5}
Dengan diagram curah hujan pada \cref{ex:g5:data:chart}: bulan mana
saja yang mendapat lebih dari $50$ mm? Berapa lebih banyak hujan
yang turun pada Mei daripada pada Juni?
\end{exercise}

\begin{exercise}[$\star$]\label{exo:g5:data:6}
Dengan diagram tanaman pada \cref{ex:g5:data:line}: berapa tinggi
tanaman itu pada minggu $2$? Di antara minggu mana saja ia tumbuh
tepat $4$ cm?
\end{exercise}

\begin{exercise}[$\star$]\label{exo:g5:data:7}
Gambarlah diagram batang untuk banyaknya gol yang dicetak sebuah
tim: September $4$, Oktober $7$, November $3$, Desember $6$.
Pilihlah skalanya, lalu sebutkan \omterm{def:g1:addition:def}{jumlahnya}.
\end{exercise}

\begin{exercise}[$\star$]\label{exo:g5:data:8}
Satu paket berisi $8$ yogurt berharga $3.20$; kalau dijual satuan,
satu yogurt berharga $0.50$. Hitunglah harga ``per satu'' di dalam
paket, dan berapa yang dihemat per yogurt dengan membeli paketnya.
\end{exercise}

\begin{exercise}[$\star\star$]\label{exo:g5:data:9}
``Sebuah bioskop menjual $148$ karcis seharga $9.50$ pada hari Sabtu
dan $95$ karcis dengan harga yang sama pada hari Minggu. Berapa
lebih banyak uang yang didapat hari Sabtu daripada hari Minggu?''
Selesaikan dengan dua cara: menghitung pemasukan kedua hari itu,
atau menghitung dengan \emph{\omterm{ex:g1:subtraction:difference}{selisih}} karcisnya --- lalu periksa
bahwa kedua jawabannya cocok.
\end{exercise}

\begin{exercise}[$\star\star$]\label{exo:g5:data:10}
Sebuah resep untuk $4$ orang memerlukan $300$ g beras. Alin memasak
untuk $10$ orang. Berapa beras yang ia perlukan? (Per satu orang
dulu --- bilangan desimal diperbolehkan.)
\end{exercise}

\begin{exercise}[$\star\star$]\label{exo:g5:data:11}
Karanglah soal tiga langkah yang memakai uang dan jawaban akhirnya
$5.50$, tulislah penyelesaian lengkapnya dengan gaya
\cref{met:g5:data:steps}, lalu ujilah pada teman sekelasmu.

\end{exercise}
